% 爱德华·康登（综述）
% license CCBYSA3
% type Wiki

本文根据 CC-BY-SA 协议转载翻译自维基百科\href{https://en.wikipedia.org/wiki/Edward_Condon}{相关文章}
\begin{figure}[ht]
\centering
\includegraphics[width=6cm]{./figures/748cfbd4f06fd931.png}
\caption{} \label{fig_ADHkd_1}
\end{figure}
爱德华·尤勒·康登（Edward Uhler Condon，1902年3月2日－1974年3月26日）是一位美国核物理学家、量子力学的先驱之一，并在第二次世界大战期间参与了雷达的研制工作，以及在曼哈顿计划中短暂参与了核武器的开发。以他命名的“弗兰克–康登原理”和“斯莱特–康登规则”均纪念了他在该领域的贡献。\(^\text{[1][2][3]}\)

他曾于1945年至1951年间担任美国国家标准局（现为美国国家标准与技术研究院，NIST）的第四任局长。1946年，康登出任美国物理学会会长；1953年，他担任美国科学促进会会长。

在麦卡锡主义时期，康登是最早成为美国众议院非美活动调查委员会攻击目标的知名科学家之一。1948年，他被公开指控为“美国原子安全体系中最薄弱的一环”，理由是他掌握大量机密信息、与原子弹研制有关联，并被指涉嫌同情共产主义和苏联。康登事件成为当时反对麦卡锡主义者（尤其是科学界）关注的标志性案件之一，舆论广泛关注，他也得到了包括美国总统哈里·杜鲁门在内的众多知名科学家的公开支持。\(^\text{[4]}\)

1968年，康登作为《康登报告》的主要作者再次广为人知。该报告是由美国空军资助的官方评估，得出的结论是“不明飞行物都有平凡的解释”。月球上的“康登环形山”即以他的名字命名，以纪念他的贡献。
\subsection{背景}
\begin{figure}[ht]
\centering
\includegraphics[width=6cm]{./figures/3223424e58a1eb6d.png}
\caption{图1. 弗兰克–康登原理能量图。由于电子跃迁相比于原子核运动要快得多，因此在跃迁过程中，最有利的振动能级是那些对应于原子核坐标变化最小的能级。图中所示的势阱表明，跃迁在 $v = 0$与$v = 2$之间最为有利。} \label{fig_ADHkd_2}
\end{figure}
爱德华·尤勒·康登于1902年3月2日出生在新墨西哥州的阿拉莫戈多，父母分别是威廉·爱德华·康登和卡罗琳·尤勒。他的父亲当时正在监督一条窄轨铁路的建设，\(^\text{[5][6]}\)当地的许多铁路都是由伐木公司修建的。1918年，康登在加利福尼亚州奥克兰完成高中学业后，曾在《奥克兰询问报》及其他报纸担任记者三年。\(^\text{[5]}\)他具有爱尔兰血统。\(^\text{[7]}\)

随后，他进入加利福尼亚大学伯克利分校就读，最初加入化学院；后来得知他高中时的物理老师成为了该校教师后，他转而改修理论物理。\(^\text{[8]}\)康登仅用三年时间获得学士学位，又用两年取得博士学位。\(^\text{[5]}\)他的博士论文结合了雷蒙德·塞耶·伯奇在带状光谱强度测量与分析方面的研究，以及詹姆斯·弗兰克的一个建议。\(^\text{[9][10]}\)

凭借国家研究委员会奖学金，康登前往德国哥廷根大学师从马克斯·玻恩，又赴慕尼黑大学师从阿诺德·索末菲。在后者指导下，康登以量子力学的形式重写了自己的博士论文，并提出了著名的“弗兰克–康登原理”。\(^\text{[9]}\)

1927年秋，他在《物理评论》上看到一则招聘广告后，加入贝尔电话实验室从事公共关系工作，主要负责宣传他们关于电子衍射的发现。\(^\text{[5][11]}\)在贝尔实验室工作期间，他还研究了词语使用频率的分布，并将结果与心理学中的韦伯–费希纳定律联系起来。\(^\text{[12]}\)
\subsection{职业生涯}
\subsubsection{早期经历}
\begin{figure}[ht]
\centering
\includegraphics[width=6cm]{./figures/6f359ea6e05a3107.png}
\caption{阿诺德·索末菲，德国理论物理学家，曾在20世纪20年代指导过康登。} \label{fig_ADHkd_3}
\end{figure}
康登曾在哥伦比亚大学短暂任教，并于1928年至1937年期间担任普林斯顿大学物理学副教授，\(^\text{[5]}\)其间曾在明尼苏达大学任教一年。\(^\text{[13]}\)他与菲利普·M·莫尔斯合著了《量子力学》，这是1929年出版的第一本英文版量子力学教材。随后，他与G.H.肖特利合著了《原子光谱理论》，该书自1935年出版起便被誉为该领域的“权威圣经”。\(^\text{[14][15][16][Note 1]}\)

1937年起，康登担任位于匹兹堡的西屋电气公司研究副主任，在那里他创建了核物理、固体物理和质谱学的研究项目。随后，他领导公司在微波雷达开发方面的研究工作。\(^\text{[13]}\)他还参与了用于分离原子弹原料铀的设备研发。\(^\text{[5]}\)在第二次世界大战期间，康登担任国家防务研究委员会的顾问，并协助组织了麻省理工学院的辐射实验室。\(^\text{[13][15]}\)1940年5月11日，康登在纽约世界博览会上展示了名为“Nimatron”的机电装置，这是一台用于玩尼姆游戏的早期计算机式设备。\(^\text{[18]}\)
\subsubsection{政府任职}
\begin{figure}[ht]
\centering
\includegraphics[width=6cm]{./figures/1ec8df323ef9ca2e.png}
\caption{J·罗伯特·奥本海默约摄于1944年在第二次世界大战期间领导了“曼哈顿计划”，而康登曾在20世纪40年代初短暂参与该计划的工作。} \label{fig_ADHkd_4}
\end{figure}
1943年，康登加入了“曼哈顿计划”。然而，仅仅六周后，他便因与项目军事负责人莱斯利·R·格罗夫斯将军在安全问题上的冲突而辞职。格罗夫斯反对康登的上级J·罗伯特·奥本海默与芝加哥大学冶金实验室主任进行讨论。[19]在辞职信中，康登表示他认为留在西屋电气公司对战争贡献更大，并解释说，洛斯阿拉莫斯实验室的安全措施令他极为不安，甚至可能使他陷入抑郁、难以有效工作。[20] 康登对奥本海默没有反抗格罗夫斯感到失望，却不知道当时奥本海默自己也尚未获得安全许可。[19]

自1943年8月至1945年2月，康登在加州大学伯克利分校兼职担任顾问，从事铀-235与铀-238的分离研究。[21] 1944年，康登当选为美国国家科学院院士。[16]1945年6月，康登受邀与多位美国著名科学家一同前往莫斯科，出席俄罗斯科学院成立220周年庆典。他表达了强烈的出席意愿。然而，当格罗夫斯得知此事后，立即联系西屋电气公司，称他认为康登此行可能存在泄露仍在进行的原子弹研究机密的风险。对此，康登试图通过白宫提出抗议。随后，格罗夫斯要求美国国务院撤销康登的护照，该请求最终被执行。[22]

战后，康登在组织科学家推动原子能由民用机构而非军事机构管理方面发挥了重要作用，[23] 主张以开放而非严格保密的方式推动原子能的发展。他担任参议员布赖恩·麦克马洪的科学顾问——麦克马洪是参议院原子能特别委员会主席，该委员会起草了《麦克马洪–道格拉斯法案》。该法案于1946年8月通过，设立了美国原子能委员会，将原子能的管理权交由民用机构。[5][15][23]康登秉持国际主义立场，主张国际科学合作，并加入了美苏科学协会。[24] 1947年，他当选为美国艺术与科学学院院士。[25]

时任美国商务部长、前副总统亨利·A·华莱士在1945年10月结识康登后，向总统推荐他出任美国国家标准局（National Bureau of Standards，现为美国国家标准与技术研究院 NIST）局长。总统哈里·S·杜鲁门接受了这一建议并提名他担任该职，参议院也一致通过了任命。康登自1945年起担任标准局局长，直至1951年离任。[26][5][21]此外，康登在1946年担任美国物理学会会长，[15][16] 并与独立公民艺术、科学与职业委员会保持联系或为其成员之一。[27] 1949年，他当选为美国哲学学会院士。[28]
\subsubsection{抨击}
\textbf{20世纪40年代}
\begin{figure}[ht]
\centering
\includegraphics[width=6cm]{./figures/da962a6675a86ce0.png}
\caption{} \label{fig_ADHkd_5}
\end{figure}
在20世纪40年代，康登的安全许可多次遭到质疑、审查，并最终被重新确认。

1946年5月29日，联邦调查局（FBI）局长约翰·埃德加·胡佛致信总统杜鲁门，信中点名多位高级政府官员涉嫌参与苏联间谍网络，并称康登“完全就是一个伪装的间谍”。（数十年后，参议员丹尼尔·帕特里克·莫伊尼汉称此指控为“毫无根据的走廊流言”。）杜鲁门政府无视了胡佛的这些指控。[29]

1947年3月21日，杜鲁门签署了美国总统行政命令第9835号，即“忠诚令”（Executive Order 9835，又称 Loyalty Order）。[30] 同月，众议员J·帕内尔·托马斯，即众议院非美活动调查委员会主席，向《华盛顿时报先驱报》提供了不利于康登忠诚度的信息，导致该报发表了两篇质疑他的报道。[31][32]托马斯选择将康登作为重点攻击目标有多重原因：他本就对科学界的国际主义精神缺乏同情，同时希望借此争议推动增加自己委员会的预算，削弱康登所支持的《麦克马洪法案》，并在选举季博取有利的媒体关注。[33]最终，美国商务部于1948年2月24日正式为康登洗清了“不忠诚”的指控。[34]
\begin{figure}[ht]
\centering
\includegraphics[width=6cm]{./figures/7e461019ca8c37fe.png}
\caption{J·帕内尔·托马斯（J. Parnell Thomas，摄于1939年）在1947年和1948年期间攻击了康登的声誉。} \label{fig_ADHkd_6}
\end{figure}
尽管如此，1948年3月2日发布的一份众议院非美活动调查委员会报告仍声称：“康登博士似乎是我国原子安全体系中最薄弱的一环。”[32][35]对此，康登回应道：“如果我真的是原子安全体系中最薄弱的一环，那倒是令人欣慰的事——因为这意味着国家可以完全放心。我是一个绝对可靠、忠诚、认真并全心全意为国家利益服务的人，这一点，我的整个生命与职业生涯都已清楚地证明。”[36]1948年3月3日，新墨西哥州民主党参议员丹尼斯·查韦斯在《国会记录》中宣读了记者马奎斯·查尔兹的一篇文章。文章指出，对康登的不利行动基于极其站不住脚的证据，而且很可能出于政治动机——尤其是与前商务部长、提名康登担任国家标准局局长的亨利·华莱士之间的政治分歧有关。[37]

1948年3月5日，来自明尼苏达州的共和党众议员乔治·麦金农在国会发言中表示：“议长先生，今日的报纸报道称，商务部长W·艾弗雷尔·哈里曼拒绝回应国会的传票，拒绝提供有关康登博士的信息。我无意对该事件的具体事实作出判断，但我必须指出，这种做法与本届政府在整个会期内对待国会传票所采取的保密态度如出一辙。”[35]1948年3月6日，《华盛顿邮报》发表社论称：“商务部长拒绝移交其部门忠诚委员会关于爱德华·U·康登博士的档案，是有充分先例可循的。”报社还反对另一项提议——将康登案的档案提交给公务员委员会下属的“忠诚审查委员会”。《邮报》指出，商务部自己的忠诚委员会已经为康登洗清了指控，该决定理应维持不变。[38]1948年3月8日，加利福尼亚州民主党众议员切斯特·E·霍利菲尔德指出：“……我想提醒各位议员注意我于1947年12月4日提交的H.R. 4641号法案。该法案旨在规范国会调查委员会的调查程序，并保障被调查人员的合法权利。如果该法案得以通过，那么像E.U.康登博士这样享誉世界的科学家，也能获得与偷鸡贼或违反交通法规者同样的权利——也就是有权为自己辩护。打着国会豁免权幌子的‘人格暗杀’，无论来自国会议员还是国会委员会，都是一种危险而可憎的嘲弄。”[35]1948年3月9日，爱达荷州民主党众议员、当时进步党副总统候选人格伦·H·泰勒，纽约州民主党众议员伊曼纽尔·塞勒，以及纽约州美国劳工党众议员利奥·伊萨克森均发表声明，谴责众议院非美活动调查委员会对康登持续不断的攻击。[35]
\begin{figure}[ht]
\centering
\includegraphics[width=6cm]{./figures/3b1ef82a1caedae1.png}
\caption{} \label{fig_ADHkd_7}
\end{figure}
康登的支持者中包括阿尔伯特·爱因斯坦和哈罗德·尤里。哈佛大学整个物理系以及众多专业组织都曾致信杜鲁门总统，为康登辩护。[39]1948年4月12日，“原子科学家紧急委员会”（举办了一场晚宴，以表明对康登的支持，活动的发起人中共有九位诺贝尔奖得主。[40]相比之下，美国国家科学院仅考虑发表一份批评众议院非美活动调查委员会调查程序的声明，而非直接为康登辩护。尽管在院士内部该声明获得了压倒性支持（275票赞成，35票反对），但科学院领导层最终决定不公开发布声明，而是选择私下与托马斯议员沟通。[41]1948年7月15日，美国原子能委员会批准了康登的安全许可，使他能够在国家标准与技术研究院接触机密信息。[42]

1948年9月，在美国科学促进会年会上，总统杜鲁门发表讲话，康登就坐于主席台附近。杜鲁门在演讲中严厉谴责托马斯议员及众议院非美活动调查委员会（HUAC），指出这种做法可能导致“在一种人人都担心被无端传言、流言蜚语和诽谤公开曝光的氛围中，重要的科学研究将无法继续进行”。他称HUAC的活动是“我们当今所面临的最不美国化的事情——这种氛围更像是一个极权国家”。[5]康登坚决反对国会试图在科学界内部进行所谓“安全风险”排查的做法。1949年6月，他在写给奥本海默的一封措辞严厉的信中批评后者向HUAC提供了有关同事的信息，信中写道：“我为此夜不能寐，怎么也想不通你怎能这样谈论一个你认识多年的朋友，一个你非常清楚他既是优秀物理学家又是好公民的人。”[19]1949年7月，康登在参议院一个就审议委员会运作规则的分委员会上作证。他批评托马斯和HUAC举行秘密听证会，并随后泄露信息以诋毁他和其他科学家的忠诚度。他指出，该委员会拒绝了他本人及其同事要求公开听证、以便自我辩护的请求。[43]
